\documentclass{beamer}
\usepackage{ctex}
\usetheme{AnnArbor}
\usecolortheme{spruce}
\usepackage[utf8]{inputenc}
\usepackage{fontspec}
\usepackage{xeCJK}
\usepackage{graphicx}
\usepackage{multicol}
\usepackage {mathtools}
\usepackage{utopia}

\setbeamerfont{headline}{family=\kaishu}
\setbeamerfont{footline}{family=\kaishu}
\setbeamerfont{frametitle}{family=\kaishu}

%----------

\definecolor{myNewColorA}{RGB}{0, 81, 40}
\definecolor{myNewColorB}{RGB}{153, 193, 173}
\definecolor{myNewColorC}{RGB}{216, 232, 224}
\setbeamercolor*{palette primary}{bg=myNewColorC}
\setbeamercolor*{palette secondary}{bg=myNewColorB, fg = white}
\setbeamercolor*{palette tertiary}{bg=myNewColorA, fg = white}
\setbeamercolor*{titlelike}{fg=myNewColorA}
\setbeamercolor*{title}{bg=myNewColorA, fg = white}
\setbeamercolor*{item}{fg=myNewColorA}
\setbeamercolor*{caption name}{fg=myNewColorA}
\setbeamercolor*{block title alerted}{bg=myNewColorA, fg=white}
\setbeamercolor*{block body alerted}{bg=myNewColorC}
\usefonttheme{professionalfonts}
\usepackage{natbib}
\usepackage{hyperref}
\hypersetup{
	colorlinks=true,
	linkcolor=black,
	citecolor=black
}

%----------

\titlegraphic{\includegraphics[height=1.5cm]{img/hzau-logo.eps}}
\setbeamertemplate{background}{\includegraphics[width=12.8cm]{img/background.jpg}}

%----------

\setbeamerfont{title}{size=\large}
\setbeamerfont{subtitle}{size=\small}
\setbeamerfont{author}{size=\small}
\setbeamerfont{date}{size=\small}
\setbeamerfont{institute}{size=\small}
\author[朱玉锋~董马丹妮~蒋运~张子栋]{朱玉锋~董马丹妮~蒋运~张子栋}
\title[Why sequence all eukaryotes?]{Why sequence all eukaryotes?}
\subtitle{为什么要测序所有真核生物？}
\institute[Group 7]{华中农业大学
	
	信息学院}
\date[2023 年 5 月 5 日]{2023 年 5 月 5 日}

%----------

\AtBeginSection[]
{
	\begin{frame}{}
		\transfade
		\tableofcontents[sectionstyle=show/shaded,subsectionstyle=show/shaded/hide]
		\addtocounter{framenumber}{0}
	\end{frame}
}

%----------

\begin{document}
	\kaishu
	
	\frame{\titlepage}

% ----------- frame -----------

	\begin{frame}
		\frametitle{分工}
	
		\begin{itemize}
			\item 董马丹妮
			\begin{itemize}
				\item 总结文章第一部分
			\end{itemize}
			
			\item 朱玉锋
			\begin{itemize}
				\item 总结文章第二部分
			\end{itemize}
			
			\item 蒋运
			\begin{itemize}
				\item 总结文章第三部分
			\end{itemize}
			
			\item 张子栋
			\begin{itemize}
				\item 材料汇总
				\item Slide 制作
			\end{itemize}
		\end{itemize}
		
	\end{frame}

% ----------- frame -----------

\begin{frame}
	\tableofcontents
\end{frame}

% ----------- frame -----------
	
	\section*{引言}
	\begin{frame}{引言}
		\begin{itemize}
			\item 作者
			\begin{itemize}
				\item 作者来自不同国家和机构的多位科学家，均为地球生物基因组计划（Earth BioGenome Project）的成员。该计划的目标时对已命名的真和生物（约 200 万种）的基因组进行高质量的测序，以尽力一个地球生命的数字图书馆。
			\end{itemize}
			\item 背景
			\begin{itemize}
				\item 文章探讨了为什么要测序所有真核生物，而不仅仅是一些代表性的物种。
			\end{itemize}
			\item 观点
			\begin{itemize}
				\item 文章认为，只有当拥有了完整的生命树或者自然生态系统中所有物种的全基因组数据时，才能回答一些关于进化和生态学的重要问题。
			\end{itemize}
		\end{itemize}
	\end{frame}

% ----------- frame -----------

	\section{发现生命之树}
	\begin{frame}{发现生命之树}
		\begin{itemize}
			\item 准确系统发育树构建的原因
			\begin{itemize}
			\item 理解表型起源及其多样化的重要数据。
			\item 杂交、多倍体化和水平基因转移在其他分支多样化中发挥重要作用。
			\item 谱系排序不完整和遗传漂移的随机性，整个基因组的基因树通常与总体物种树不同。
			\end{itemize}

			\item 准确系统发育树构建的难题
			\begin{itemize}
			\item 国际核苷酸数据库中物种序列少，物种之间同源基因重叠低。
			\item 草稿形式基因组使用范围小，单细胞生物深层谱系位置未解决。
			\end{itemize}
		\end{itemize}

	\end{frame}


% ----------- frame -----------

	\begin{frame}{发现生命之树}
		\begin{itemize}
		\item 可以解决的问题
		\begin{enumerate}
			\item 杂交和多倍体化在生命树中的普遍性、重要性；基因组中留下的古老的杂交、多倍体化和渐渗特征，
			\item 新基因具体产生机制，不同谱系中可预测的基因家族起源和多样化模式的存在性。
			\item 水平基因转移在真核生物进化中的重要性。
		\end{enumerate}
		\item 意义
		\begin{enumerate}
			\item 为「真核生物树所有分支的起源和多样化的产生」这一假设提供充分支持，绘制网状事件和基因复制、分化和丢失的时间。
			\item 完整的序列更可能包含与感兴趣的节点、过程或事件有关的信息。
			\item 为不同问题选择合适的数据。
			\item 突出功能 DNA、RNA 和氨基酸残基的差异和限制，为生物工程提供新功能的信息。
		\end{enumerate}
	\end{itemize}
	\end{frame}

% ----------- frame -----------

	\section{定义真核细胞的起源}
	\subsection*{需要解决的问题}
	\begin{frame}{定义真核细胞的起源}{需要解决的问题}
		\begin{enumerate}
			\item 对于真核细胞复杂性的前因以及所涉及的环境和选择性因素知之甚少。
			\item 目前远不清楚参与以下过程中的单个蛋白同源物如何、何时、以何种顺序进化：
			\begin{itemize}
				\item 细胞核
				\item 时空分离的转录和翻译
				\item 线粒体
				\item 线性染色体
				\item 减数分裂
				\item 有丝分裂
				\item 内膜和蛋白质运输系统以及鞭毛
			\end{itemize}
		\end{enumerate}

	\end{frame}

% ----------- frame -----------
	
\begin{frame}{定义真核细胞的起源}{完整真核基因组可以解决的问题}
	\begin{itemize}
		\item 所有现存真核生物的祖先中都存在哪些基因？
		\item 真核细胞的复杂性是以逐步的方式进化还是一次进化？
		\item 真核细胞的复杂性是以逐步的方式进化还是一次进化？
		\item 哪些基因组过程限制和促进了次生和次生内共生生物光合作用的水平传播？
	\end{itemize}

	猜测：对真核细胞起源的变革性见解很可能来自古菌和微生物真核结构域内目前未知的细胞多样性序列。高分辨率真核生物生命树的生成将为核心真核生物工具包和器官功能中基因的起源和多样性分析奠定基础。

	\begin{alertblock}{总结}
		一套完整的真核基因组可以解决关于真核细胞进化起源的基本问题。
	\end{alertblock}
\end{frame}

% ----------- frame -----------
	
	\section{追踪共生的基因组变化}
	\begin{frame}{追踪共生的基因组变化}

		\begin{itemize}
			\item 共生的互惠关系
			\begin{enumerate}
				\item 维管植物和定根微生物之间的互惠共生关系与昆虫营养生态位。
				\item 共生体之间再面对免疫和反编辑系统施加的障碍时存在依赖关系。
				\item 物种存在于有助于食物消化和其他过程的微生物组中，微生物组的基因组可能与理解宿主生理学有关。
			\end{enumerate}
		\end{itemize}

		\begin{alertblock}{需要解决的问题}
			共生对真核细胞生态多样性和生物复杂性的影响。
		\end{alertblock}
	\end{frame}

% ----------- frame -----------

	\begin{frame}{追踪共生的基因组变化}
		\begin{itemize}
			\item 可以解决的问题
			\begin{itemize}
			\item 宿主及其微生物组是否视为全基因组超有机体。
			\item 寄生共生体和互惠共生体基因组的进化是否有共同的主题。
			\item 在基因和分子水平上实现共生机制的范围和多样性。
			\item 有益但非强制性的共生体对双方进化的影响。
			\item 共生基因转移的普遍性，共生体消失的遗产对当今生物功能的贡献。
			\item 大型宿主生物及其微生物群之间的共同进化对多样化的推进。
			\end{itemize}

			\item 作用
			\begin{enumerate}
				\item 发现新的基因和调节机制。
				\item 共生体之间相对进化轨迹有助于了解这些生物之间的必要联系。
				\item 了解害虫和寄生虫及其互惠共生体的基本生理依赖性。
			\end{enumerate}
			
		\end{itemize}
	\end{frame}

% ----------- frame -----------

	\section{解密染色体进化}
	\begin{frame}{解密染色体进化}
		\begin{itemize}
			\item 多数真核生物的基因组是以线性染色体组织的二倍体。
			\item 真核染色质的三维构象对基因表达和细胞分化至关重要。
			\item 真核生物通过添加各种非基因组模板DNA，保护其染色体末端端粒。
			\item 基因在不同染色体上的分配为基因调控网络进化提供机会和约束。
			\item 多数真核生物类群中，核型和同源性通常保守。
		\end{itemize}
	\end{frame}

% ----------- frame -----------

	\begin{frame}{解密染色体进化}
		\begin{itemize}
			\item 真核基因组可解决的问题
			\begin{itemize}
				\item 基因组的可能核型、3D组织和同源关系。
				\item 新物种的形成过程中，染色体数量稳定方式。
				\item 有丝分裂和减数分裂对染色体数量和组织的限制。
				\item 染色体重排是否为真核生物进化中转变和适应的原因。
				\item 端粒进化方式及不同端粒维持机制的意义。
				\item 着丝粒种类对基因组进化的影响。
				\item 核型对基因调控网络进化的限制。
			\end{itemize}

			\item 作用
			\begin{enumerate}
				\item 可解决所有规模的染色体组织进化问题，从局部同源性的保守性分析到核型进化，以及着丝粒和端粒重新定位的动力学。
				\item 可以支持开发具有工程染色体稳定性和其他行为的新型合成基因组。
			\end{enumerate}
		\end{itemize}
	\end{frame}

% ----------- frame -----------

	\section{揭示真核基因调控的深层逻辑}
	\begin{frame}{揭示真核基因调控的深层逻辑}
		\begin{multicols}{2}
			\begin{enumerate}
				\item 每个真核基因组都包含数千个蛋白质编码基因，每个基因都有调控元件和回路。这些元件中的大多数可以与非功能DNA区分开来。
				\item 特定谱系变化速度加快的区域可能与新的进化适应有关。对于单个物种，参考基因组可以测定组织和细胞中的调节。
				\item 可以将在一个物种中实验确定的功能注释映射到相关物种，从而提供整个分支的快速、深入的基因组注释。
			\end{enumerate}
			\begin{itemize}
				\item 可解决的问题
				\begin{itemize}
				\item 所有真核基因的综合列表。
				\item 真核生物基因组调控的组成部分及其稳定性。
				\item 新的监管制度的产生、对取代现有制度的取代。
				\item 深层多物种全基因组比对能否揭示新的保守元素类别。
				\item 哪些基因组特征区分保守的调控网络。
				\item 调控网络限制染色体和基因组进化的方式。
				\end{itemize}
			\end{itemize}
		\end{multicols}
	\end{frame}

% ----------- frame -----------

	\begin{frame}{揭示真核基因调控的深层逻辑}
		\begin{itemize}
			\item 作用
			\begin{enumerate}
				\item 可以以更高的灵敏度和更高的分辨率绘制功能序列，并将功能变异与表型联系起来。
				\item 可以绘制长程染色体相互作用图，并解开重复、移动元件插入和重复积累在修改先前存在的基因调控中的影响。
			\end{enumerate}

			\quad

			\item 意义
			\begin{enumerate}
				\item 将有可能绘制蛋白质编码基因和长非编码RNA的调控启动子区、重复、重复和转座元件插入，是基因调控创新的关键来源。
				\item 将每个物种单独考虑，可以更深入地了解调控基因组及其功能。
			\end{enumerate}
		\end{itemize}
	\end{frame}

% ----------- frame -----------

	\section{探索性系统的多样性}
	\begin{frame}{探索性系统的多样性}
		\begin{columns}
			\column{0.5\textwidth}
			\begin{itemize}
				\item 性系统的多样性
				\begin{enumerate}
					\item 一些原生生物群体（纤毛虫）有多种等效的交配类型，但大多数多细胞生物有两性。
					\item 分化单倍体配子（大卵子和小精子）的产生多次独立进化，父母对卵子或精子的贡献导致了形态、生活史策略和行为上极端性别差异的进化。
					\item 决定性别的机制多种多样，从单基因座驱动因素和分化的性染色体到表观遗传和环境决定。
				\end{enumerate}
			\end{itemize}

			\column{0.5\textwidth}
			\begin{itemize}
				\item 真核谱系完整测序可以解决的问题
				\begin{itemize}
					\item 雄性-雌性系统多久独立进化一次？具有多种交配类型的系统如何进化和维持？
					\item 性别决定系统的变化会驱动物种形成吗？如何在基因组中反映？
					\item 在具有分离种系的物种（纤毛虫和大多数动物）中，种系如何指定？种系基因组如何维持？
					\item 染色质减少如何在生命树中分布及调节的？
					
				\end{itemize}
			\end{itemize}
		\end{columns}
	\end{frame}

% ----------- frame -----------

	\section{探索物种基因组学中的多样性}
	\begin{frame}{探索物种基因组学中的多样性}
		\begin{itemize}
			\item 争议
			\begin{enumerate}
				\item 物种形成的过程。
				\item 许多情况下物种边界是泄漏或不完整的。
				\begin{itemize}
					\item 生物物种概念认为姐妹分类群之间存在稳定的繁殖隔离。然而，基因组数据表明，在分类学上被接受的物种形成后，通过杂交和渗入进行的基因交换是常见的。
				\end{itemize}
			\end{enumerate}
			\quad
			\item 探索物种基因组学多样性的原因
			\begin{itemize}
				\item 命名物种的概念隐藏了多个神秘的分类群，许多物种尚未被发现和命名。基因组数据对发现和定义新的分类群至关重要。
			\end{itemize}
		\end{itemize}
	\end{frame}

% ----------- frame -----------

	\subsection*{真核生物谱系完整测序可解决物种形成问题}
	\begin{frame}{探索物种基因组学中的多样性}{真核生物谱系完整测序可解决物种形成问题}
		\begin{itemize}
			\item 物种形成过程的多样性以及它们是如何形成基因组的？
			\item 哪些基因组信号区分不完全谱系分类、渗入和杂交？
			\item 某些类型的基因和基因网络更有可能与物种形成有关吗？
			\item 染色体重排在物种形成中扮演什么角色？
			\item 不同的繁殖策略是否会驱动不同的基因组结构，从而影响物种形成？
			\item 物种辐射等宏观进化现象是如何塑造基因组的？
			\item 从目前的基因组来看，未来的物种形成是可以预测的吗？
		\end{itemize}
	\end{frame}

% ----------- frame -----------

	\section*{Q\&A}
	\begin{frame}
		\begin{center}
			\textcolor{myNewColorA}{\Huge {Q\&A}}
		\end{center}
	\end{frame}

% ----------- frame -----------

	\section*{致谢}  
	\begin{frame}
		\begin{center}
			\textcolor{myNewColorA}{\Huge {谢谢!\\ \quad \\ 请老师批评指正！}}
		\end{center}
	\end{frame}
	
\end{document}